% Generated draft: requires mathematical and source review.
\documentclass[11pt]{article}
\usepackage[margin=25mm]{geometry}
\usepackage{fontspec}
\setmainfont{Latin Modern Roman}
\usepackage{amsmath,amssymb,mathtools,unicode-math}
\setmathfont{Latin Modern Math}
\usepackage{enumitem,needspace,xurl,hyperref}
\setlength{\parindent}{0pt}
\setlength{\parskip}{5pt}
\setlength{\emergencystretch}{3em}
\newcommand{\N}{\mathbb{N}}
\newcommand{\Z}{\mathbb{Z}}
\newcommand{\Q}{\mathbb{Q}}
\newcommand{\R}{\mathbb{R}}
\newcommand{\C}{\mathbb{C}}
\newcommand{\F}{\mathbb{F}}
\newcommand{\A}{\mathbb{A}}
\newcommand{\PP}{\mathbb{P}}
\newcommand{\E}{\mathbb{E}}
\newcommand{\eps}{\varepsilon}
\newcommand{\dd}{\,\mathrm{d}}
\begin{document}
\section*{UnsolvedMath ID: 30003835}
\textbf{Source problem code:} catalogue entry
\subsection{Definitions and mathematical statement}
Let \(M\) be a connected, oriented one-dimensional manifold, and let \(\operatorname{Homeo}^{+}(M)\) denote its group of orientation-preserving homeomorphisms. Let \(G<\operatorname{Homeo}^{+}(M)\), and denote the given faithful action by \(\iota:G\hookrightarrow\operatorname{Homeo}^{+}(M)\). The source does not specify whether manifolds with boundary are allowed.\par
For a specified differentiability class \(C^r\) and a compatible differentiable structure on \(M\), write \(\operatorname{Diff}^{r}(M)\) for the group of bijections of \(M\) whose maps and inverses are of class \(C^r\). The source leaves the range of \(r\) and the differentiable structure unspecified; the following questions are conditional on making these choices.
\begin{enumerate}
\item \emph{Abstract smoothability.} For the given \(G\) and \(M\), does there exist an injective homomorphism
\[
\rho:G\longrightarrow\operatorname{Diff}^{r}(M)?
\]
Here \(G\) is considered as an abstract group, so this question imposes no continuity requirement on \(\rho\) with respect to a topology inherited from \(\operatorname{Homeo}^{+}(M)\). The source does not explicitly require this new action to preserve orientation.
\item \emph{Smoothability of the given action.} Does there exist a homeomorphism \(h:M\to M\) such that
\[
h\,\iota(g)\,h^{-1}\in\operatorname{Diff}^{r}(M)
\qquad\text{for every }g\in G?
\]
This asks for one homeomorphism simultaneously conjugating the entire given action to a \(C^r\) action.
\item \emph{Obstructions.} What properties of \(G\), of the given action, of \(M\), or of the regularity \(r\) obstruct the actions or conjugacies above? This is a request for obstruction criteria, with no specific necessary or sufficient condition asserted by the source.
\end{enumerate}
The record also asks whether an infinite Kazhdan group can act faithfully by homeomorphisms on the circle \(S^1=\R/\Z\). A Kazhdan group is a group with property \((T)\). In the discrete-group interpretation, property \((T)\) means that every unitary representation \(\pi:K\to\mathcal U(H)\) on a complex Hilbert space having almost invariant unit vectors has a nonzero invariant vector. Here almost invariant unit vectors means
\[
\forall E\subseteq K\text{ finite}\ \forall\varepsilon>0\ \exists v\in H:
\quad \|v\|=1,\qquad
\|\pi(k)v-v\|<\varepsilon\quad(k\in E).
\]
Under an existential reading, the circle question is
\[
\exists K\ \bigl[\,|K|=\infty,\ K\text{ has property }(T),\ 
\exists\text{ an injective homomorphism }
K\to\operatorname{Homeo}(S^1)\,\bigr]?
\]
The wording does not explicitly distinguish this reading from asking whether every infinite Kazhdan group admits such an action. Moreover, the statement does not impose countability, whereas the supplied literature assessment specifically mentions faithful \emph{countable} property-\((T)\) circle actions. These quantifier and domain qualifications remain unresolved. The circle question uses all homeomorphisms; orientation preservation is not expressly required.
\subsection{Short English statement}
For a group already acting faithfully on a connected one-dimensional manifold, ask whether its abstract group has a faithful \(C^r\) action on the same manifold, whether the given action becomes \(C^r\) after a topological change of coordinates, and what obstructs these possibilities. A further question concerns faithful circle actions of infinite groups with Kazhdan's property \((T)\). Regularity, boundary conventions, and the precise quantifiers and countability restriction in the Kazhdan question require clarification.
\subsection{Sources}
\begin{itemize}
\item Supplied local UnsolvedMath registry record: permanent ID \(30003835\), title \emph{Smoothability of Group Actions on One-Manifolds}, exact source problem code \texttt{catalogue entry}, in the Oberwolfach Reports Open Problems collection. The supplied proposed year is 2018; no individual proposer is named.
\item Supplied report citation: \emph{Cohomological and Metric Properties of Groups of Homeomorphisms of R} (2019), Oberwolfach Reports, pp. 1579--1633, DOI \texttt{10.4171/owr/2018/26}. \url{https://doi.org/10.4171/owr/2018/26}. The background attributes the collection to workshop reports and their open-problem discussions; no exact problem page or named speaker is supplied.
\item Supplied literature reference: Lodha--Matte Bon--Triestino, \url{https://doi.org/10.1112/topo.12151}.
\item Supplied literature reference: Duchesne, \url{https://arxiv.org/abs/2011.12861}.
\item Supplied literature reference: Guba, \url{https://doi.org/10.4171/JCA/110}.
\item Literature assessment dated exactly \texttt{2026-08-21}: the supplied curation labels the bundle open, reports partial smoothability and rigidity theorems, and states that faithful countable property-\((T)\) circle actions and amenability of \(F\) remain unresolved. This is supplied dated curation, not an independently verified literature assessment. The background expressly requests independent verification of time-sensitive status and cited claims. The mention of amenability of \(F\) is preserved as background; the source statement neither defines \(F\) nor includes its amenability as a formulated question.
\item The registry supplies the status \texttt{open} and the statement-status label \texttt{corrected\_verified}, accompanied by a claim of manual semantic, typographic, and LaTeX review. These are supplied registry labels and review claims, not independent verification of this draft or of the mathematical status.
\end{itemize}
\end{document}
